\documentclass[11pt]{article}

\usepackage[utf8]{inputenc}
\usepackage[T1]{fontenc}
\usepackage{lmodern}
\usepackage{geometry}
\usepackage{amsmath,amssymb,amsfonts,amsthm,mathtools}
\usepackage{bm}
\usepackage{enumitem}
\usepackage{microtype}
\usepackage{hyperref}
\usepackage{titlesec}
\usepackage{booktabs}
\usepackage{array}
\usepackage{setspace}
\usepackage{mathrsfs}
\usepackage{physics}

\geometry{margin=1in}
\hypersetup{colorlinks=true, linkcolor=black, urlcolor=blue, citecolor=black}
\setlist[enumerate]{topsep=0.4em,itemsep=0.35em}
\setlist[itemize]{topsep=0.3em,itemsep=0.25em}
\titleformat{\section}{\large\bfseries}{\thesection.}{0.5em}{}
\titleformat{\subsection}{\normalsize\bfseries}{\thesubsection.}{0.5em}{}
\setstretch{1.06}

\newcommand{\Aether}{\AE{}ther}
\newcommand{\Aflow}{\AE{}ther-flow}
\newcommand{\TheoryName}{\Aflow{} Interpretation of Relativity}
\newcommand{\TheoryProgram}{\Aether{} / \Aflow{} framework}
\newcommand{\Sub}{\mathcal{A}}
\newcommand{\BridgeMetric}{\mathfrak{g}}
\newcommand{\ReBridge}{\widetilde{\mathfrak{g}}}
\newcommand{\Proj}{\mathcal D}
\newcommand{\Dperp}{\mathcal D}
\newcommand{\PiTheta}{\Pi^{(\vartheta)}}
\newcommand{\ReEff}{\widetilde{\mathcal T}_{\mu\nu}^{\mathrm{br,eff}}}
\newcommand{\ReLong}{\widetilde{\mathcal R}_{\mu\nu}^{(\chi,\ge 3)}}
\newcommand{\ReOff}{\widetilde{\mathcal R}_{\mu\nu}^{\mathrm{off}}}
\newcommand{\AY}{\mathcal A_Y}
\newcommand{\AZ}{\mathcal A_Z}

\newtheorem{definition}{Definition}
\newtheorem{proposition}{Proposition}
\newtheorem{remark}{Remark}
\newtheorem{corollary}{Corollary}
\newtheorem{theorem}{Theorem}

\title{\textbf{The \TheoryName{}}\\[0.4em]
\large Redesigned Bridge-Closure Reduction and Tensor Verdict}
\author{}
\date{}

\begin{document}

\maketitle

\begin{abstract}
The post-quotient bridge-metric redesign note proved only a narrow same-scope fact. With the revised bridge
\[
\ReBridge_{AB}
=
G_{AB}-2U_AU_B-\frac{1}{\kappa_\theta}\PiTheta_{AB},
\]
and with matter still coupled through the single operational metric
\[
g^{\mathrm{br}}_{\mu\nu}=\Xi^\ast\ReBridge_{AB},
\]
the old longitudinal quotient remainder
\[
T_{\mu\nu}^{(\chi,\mathrm{eff},2)}
=
\frac{1}{\kappa_\theta}
\left(
\Dperp_\mu\Phi_K\,\Dperp_\nu\Phi_K
-\frac12 g_{\mu\nu}\Dperp_\alpha\Phi_K\,\Dperp^\alpha\Phi_K
\right)
\]
is cancelled at that same tensor-reduction scope. The present note performs the next, broader test that was still missing: derive the full eliminated observer equation on the revised bridge, check whether any higher-order or off-longitudinal remainder survives, and decide whether the minimal convergence-response bridge really closes the observer tensor.

The result is negative for full closure. On reduced redesigned solutions one has
\[
G_{\mu\nu}[g^{\mathrm{br}}]
=
8\pi G_N\,T_{\mu\nu}^{\mathrm{matter}}
+ \ReEff,
\]
and after the same slice elimination already used on the quotient branch, the redesigned effective tensor splits as
\[
\ReEff
=
T_{\mu\nu}^{(Q,\mathrm{br,eff})}
+T_{\mu\nu}^{(W,\mathrm{br,eff})}
+\ReLong.
\]
Here \(\ReLong=\mathcal O_{\ge 3}\) is the surviving higher-order longitudinal remainder left after the quadratic same-scope cancellation, while the off-longitudinal piece
\[
\ReOff
\equiv
T_{\mu\nu}^{(Q,\mathrm{br,eff})}
+T_{\mu\nu}^{(W,\mathrm{br,eff})}
\]
is untouched at leading order because the bridge redesign depends only on the convergence potential \(\Phi_\vartheta\).

The redesigned bridge therefore fails the full closure test already on a convergence-free rotational regime. If \(K\equiv0\) but the transverse vorticity solve is nontrivial, then \(\Phi_\vartheta\equiv0\), so the new bridge slot vanishes, yet
\[
\Gamma_{W,\mathrm{eff}}^{\mathrm{br},(2)}
=
-\frac{\kappa_\Omega}{4}
\int_{\mathbb R^{1+3}}\Omega^T_{ij}\Omega^{T,ij}\,dt\,d^3x
\]
produces the nonzero tensor
\[
T_{\mu\nu}^{(W,\mathrm{br,eff},2)}
=
\kappa_\Omega
\left(
\Omega_{\mu\alpha}^T\Omega_\nu^{T\,\alpha}
-\frac14 g_{\mu\nu}^{\mathrm{br}}\Omega_{\alpha\beta}^T\Omega^{T,\alpha\beta}
\right),
\qquad
\rho_{W}^{\mathrm{br},(2)}
=
\frac{\kappa_\Omega}{4}\abs{\Omega^T}^2.
\]
Hence the minimal convergence-response bridge does not close the observer tensor. It removes the old same-scope longitudinal slot but leaves a genuine redesigned observer remainder, both through the unresolved higher-order longitudinal sector and, decisively, through off-longitudinal auxiliary stress. The verdict is that the first bridge-first redesign only postpones the problem. With this full redesigned-bridge test now negative, it is finally honest to open the second redesign note that reconsiders the single-metric matter route.
\end{abstract}

\section{Introduction}

The positive quotient line remains positive where it was already positive. The ordered-flow-label quotient candidate, the separate quotient-architecture screen, the quotient reduced-system note, the quotient local/coercive note, the quotient theorem note, and the quotient decay and asymptotic notes still give a clean dynamical branch with
\[
\AY=\AZ=0
\]
imposed from the outset and with no physical or analytic \(S,Y,Z\) data reintroduced. None of that work is reopened here.

What changed after the negative quotient tensor verdict was narrower. The fixed bridge metric
\[
\BridgeMetric_{AB}=G_{AB}-2U_AU_B
\]
left a genuine eliminated observer remainder, and the first post-verdict redesign note therefore tested the smallest bridge-only move: add one convergence-response tensor \(\PiTheta_{AB}\), keep the single-metric matter route fixed, and ask whether the old longitudinal quotient remainder disappears at the same tensor-reduction scope. That same-scope test was positive.

The same-scope test did \emph{not} prove full closure. It did not recompute the actual eliminated observer equation on the revised bridge. It did not decide whether higher-order longitudinal terms survive. It did not decide whether off-longitudinal \(Q\)- or vorticity-sector remainders survive. The present note performs exactly that missing step.

\paragraph{Framework and claim boundary.}
\TheoryProgram{} denotes the broader \Aether{} / \Aflow{} framework built on the claims that the \Aether{} is the underlying four-dimensional substrate of reality, the \Aflow{} is its intrinsic ordered motion, observed three-dimensional space is the local experiential slice of that deeper substrate, S-time is the experienced order of change arising from matter, light, and the \Aflow{}, and observed expansion is the three-dimensional appearance of deeper four-dimensional ordered motion. Within that framework, \TheoryName{} denotes the exact-closure theory in which the effective gravitational dynamics are adopted to be exactly Einsteinian with universal matter coupling. In the active sequence, \emph{adoption} means use of that established relativistic dynamics without claiming substrate derivation, while \emph{derivation} is reserved for a first-principles recovery from explicit substrate variables.

The present manuscript stays strictly inside that boundary. It does not claim a completed substrate derivation. It does not reopen reduced-system, theorem, decay, or asymptotic work. It keeps the positive quotient branch and the single-metric matter route fixed, computes the actual eliminated observer equation for the revised bridge metric, and records the tensor-level verdict forced by that computation.

\section{Current Redesigned Branch}

\begin{definition}[Current redesigned bridge branch]
\label{def:branch}
The \emph{current redesigned bridge branch} consists of:
\begin{enumerate}
    \item the same reduced quotient action and the same quotient equations
    \begin{equation}
    \AY=0,
    \qquad
    \AZ=0,
    \label{eq:AYAZ}
    \end{equation}
    already used on the positive quotient line;
    \item the revised bridge metric
    \begin{equation}
    \ReBridge_{AB}
    =
    G_{AB}-2U_AU_B-\frac{1}{\kappa_\theta}\PiTheta_{AB},
    \label{eq:newbridge}
    \end{equation}
    with
    \begin{equation}
    P_{AB}=G_{AB}-U_AU_B,
    \qquad
    \vartheta=\nabla_AU^A,
    \qquad
    \Proj_A=P_A{}^B\nabla_B,
    \label{eq:proj}
    \end{equation}
    and
    \begin{equation}
    \PiTheta_{AB}
    =
    \Proj_A\Phi_\vartheta\,\Proj_B\Phi_\vartheta
    -\frac12 P_{AB}\Proj_C\Phi_\vartheta\,\Proj^C\Phi_\vartheta,
    \qquad
    (-\Delta_\perp)\Phi_\vartheta=\vartheta;
    \label{eq:pitheta}
    \end{equation}
    \item the revised operational metric
    \begin{equation}
    g^{\mathrm{br}}_{\mu\nu}=\Xi^\ast\ReBridge_{AB};
    \label{eq:gbr}
    \end{equation}
    \item the same single-metric matter route
    \begin{equation}
    S_{\mathrm{matter}}^{\mathrm{br}}[G,U,\Xi,\psi]
    =
    S_{\mathrm{matter}}[\Xi^\ast\ReBridge,\psi]
    =
    S_{\mathrm{matter}}[g^{\mathrm{br}},\psi].
    \label{eq:matter}
    \end{equation}
\end{enumerate}
\end{definition}

\begin{remark}
Definition \ref{def:branch} changes only the bridge metric. The quotient physical field space, the quotient slice package, and the already-positive reduced-system/theorem/decay/asymptotic line are not modified structurally.
\end{remark}

\begin{proposition}[The slice package stays of the same type]
\label{prop:slice}
On the redesigned bridge branch, the slice-wise nonpropagating variables still satisfy an elliptic package of the same principal type:
\begin{align}
\bigl((-\Delta_\gamma)\delta_{IJ}+\widehat M_{IJ}\bigr)Q^J
&=
\mathcal N_I^{(Q,\mathrm{br})},
\label{eq:Qeq}
\\
\kappa_\theta\Delta_\gamma\chi
&=
-K+\mathcal N_\chi^{\mathrm{br}},
\label{eq:chieq}
\\
\kappa_\Omega(-\Delta_\gamma)W_i^T
&=
\mathcal N_i^{(T,\mathrm{br})},
\label{eq:Weq}
\end{align}
with \(\AY=\AZ=0\) and with modified only lower-order source terms.
\end{proposition}

\begin{proof}
The redesign note already established that the bridge term \(\PiTheta_{AB}\) does not annihilate the ordered-flow equation and does not introduce a new propagating scalar channel. The positive quotient branch still carries the same \(Q\)-, divergence-, and vorticity-sector auxiliary blocks. Hence the principal operators of the slice package remain exactly the positive projected-mass \(Q\)-operator, the longitudinal \(\kappa_\theta\Delta_\gamma\) operator, and the transverse \(\kappa_\Omega(-\Delta_\gamma)\) operator. Only lower-order couplings are altered by the revised bridge metric.
\end{proof}

\section{Full Eliminated Observer Equation on the Revised Bridge}

\begin{definition}[Redesigned eliminated auxiliary action]
\label{def:Gamma}
Let
\[
\widetilde{\mathscr S}_{\mathrm{br}}[g^{\mathrm{br}},\psi]
=
\bigl(Q[g^{\mathrm{br}},\psi],\chi[g^{\mathrm{br}},\psi],W^T[g^{\mathrm{br}},\psi]\bigr)
\]
denote the decaying slice solver for \eqref{eq:Qeq}--\eqref{eq:Weq}. The \emph{redesigned eliminated auxiliary action} is
\begin{equation}
\widetilde{\Gamma}_{\mathrm{aux}}^{\mathrm{br}}[g^{\mathrm{br}},\psi]
\equiv
\int_{\Sub} d^4X\,\sqrt{G}\,
\mathcal L_{\mathrm{aux,deg}}
\Big|_{\widetilde{\mathscr S}_{\mathrm{br}}[g^{\mathrm{br}},\psi]},
\label{eq:Gamma}
\end{equation}
where \(\mathcal L_{\mathrm{aux,deg}}\) is the same reduced \(S\)-free auxiliary Lagrangian already used on the quotient branch.
\end{definition}

\begin{definition}[Redesigned effective observer tensor]
\label{def:Teff}
Define the \emph{redesigned effective observer tensor} by
\begin{equation}
\ReEff
\equiv
-\frac{2}{\sqrt{-g^{\mathrm{br}}}}
\frac{\delta\widetilde{\Gamma}_{\mathrm{aux}}^{\mathrm{br}}}{\delta (g^{\mathrm{br}})^{\mu\nu}}.
\label{eq:Teff}
\end{equation}
\end{definition}

\begin{proposition}[Actual eliminated observer equation]
\label{prop:redeq}
On reduced redesigned solutions satisfying the slice equations and \eqref{eq:AYAZ},
\begin{equation}
G_{\mu\nu}[g^{\mathrm{br}}]
=
8\pi G_N\,T_{\mu\nu}^{\mathrm{matter}}
+\ReEff.
\label{eq:redeq}
\end{equation}
\end{proposition}

\begin{proof}
The observer metric has been changed from \(\BridgeMetric\) to \(\ReBridge\), but the action still has the same three pieces: the Einstein--Hilbert term of the operational bridge metric, the reduced auxiliary action, and the matter action coupled through that same operational metric. As in the earlier quotient bridge-closure reduction note, substituting the slice solver and varying with respect to the operational metric produces chain-rule terms proportional to the slice equations, which vanish on reduced solutions. What remains is the metric variation of the eliminated auxiliary action, namely \eqref{eq:Teff}. This yields \eqref{eq:redeq}.
\end{proof}

\section{Longitudinal Cancellation Versus Higher-Order Longitudinal Remainder}

The previous redesign note computed only the first quadratic longitudinal slot. The first step in the full closure test is to isolate what remains of that sector after the same-scope cancellation has been imposed.

\begin{proposition}[Quadratic longitudinal cancellation is exact only at that slot]
\label{prop:longcancel}
On the redesigned branch with coefficient \(-1/\kappa_\theta\), the longitudinal contribution to \(\ReEff\) has the form
\begin{equation}
T_{\mu\nu}^{(\chi,\mathrm{br,eff})}
=
\ReLong,
\qquad
\ReLong=\mathcal O_{\ge 3},
\label{eq:longres}
\end{equation}
where \(\ReLong\) is the residual higher-order longitudinal tensor left after the cancellation of the quadratic slot
\[
\frac{1}{\kappa_\theta}
\left(
\Dperp_\mu\Phi_K\,\Dperp_\nu\Phi_K
-\frac12 g_{\mu\nu}\Dperp_\alpha\Phi_K\,\Dperp^\alpha\Phi_K
\right).
\]
\end{proposition}

\begin{proof}
The post-quotient redesign note already proved that
\[
\Xi^\ast\PiTheta_{\mu\nu}
=
\Dperp_\mu\Phi_K\,\Dperp_\nu\Phi_K
-\frac12 g_{\mu\nu}\Dperp_\alpha\Phi_K\,\Dperp^\alpha\Phi_K
+\mathcal O_{\ge 3},
\]
so the coefficient choice \(-1/\kappa_\theta\) removes exactly the quadratic longitudinal tensor computed in the previous quotient tensor-verdict note. All remaining longitudinal terms arise from the nonlinear difference between \(\vartheta\) and \(K\), the nonlinear part of the slice source \(\mathcal N_\chi^{\mathrm{br}}\), and the nonlinear metric dependence of the same eliminated longitudinal functional. Those terms all start at cubic order, giving \eqref{eq:longres}.
\end{proof}

\begin{remark}
Proposition \ref{prop:longcancel} is the precise meaning of the earlier same-scope positive result. The minimal convergence-response bridge clears the old quadratic longitudinal slot, but it does not prove any all-order longitudinal identity.
\end{remark}

\section{Off-Longitudinal Effective Tensor}

The full closure question is broader than the longitudinal slot alone. The redesigned bridge must also be tested against the off-longitudinal eliminated tensor carried by the \(Q\)- and vorticity-sector auxiliaries.

\begin{definition}[Off-longitudinal redesigned remainder]
\label{def:off}
Define the \emph{off-longitudinal redesigned remainder} by
\begin{equation}
\ReOff
\equiv
T_{\mu\nu}^{(Q,\mathrm{br,eff})}
+T_{\mu\nu}^{(W,\mathrm{br,eff})},
\label{eq:off}
\end{equation}
where \(T_{\mu\nu}^{(Q,\mathrm{br,eff})}\) and \(T_{\mu\nu}^{(W,\mathrm{br,eff})}\) are the contributions to \eqref{eq:Teff} coming from the eliminated \(Q\)- and transverse-vorticity blocks.
\end{definition}

\begin{proposition}[Full redesigned tensor decomposition]
\label{prop:decomp}
On reduced redesigned solutions,
\begin{equation}
\ReEff
=
\ReOff+\ReLong.
\label{eq:decomp}
\end{equation}
\end{proposition}

\begin{proof}
The reduced auxiliary action splits into its \(Q\)-sector, longitudinal, and transverse-vorticity pieces exactly as in the earlier quotient tensor-verdict note. Proposition \ref{prop:longcancel} replaces the old explicit quadratic longitudinal tensor by the higher-order residual \(\ReLong\). The remaining sectors are precisely the off-longitudinal \(Q\)- and \(W\)-pieces. This yields \eqref{eq:decomp}.
\end{proof}

\begin{proposition}[The minimal convergence-response bridge is blind to the off-longitudinal sectors]
\label{prop:blind}
At leading order, the revised bridge term \(\PiTheta_{AB}\) does not alter the off-longitudinal effective tensor. More precisely,
\begin{equation}
T_{\mu\nu}^{(Q,\mathrm{br,eff})}
=
T_{\mu\nu}^{(Q,\mathrm{eff})}
+\mathcal O_{\ge 3},
\qquad
T_{\mu\nu}^{(W,\mathrm{br,eff})}
=
T_{\mu\nu}^{(W,\mathrm{eff})}
+\mathcal O_{\ge 3}.
\label{eq:blind}
\end{equation}
\end{proposition}

\begin{proof}
The redesigned bridge term depends only on the convergence potential \(\Phi_\vartheta\). It introduces no dependence on the projected spatial \(Q\)-block and no new dependence on the transverse vorticity variable \(W_i^T\) beyond the background metric factors already present in the quotient branch. Therefore the \(Q\)- and \(W\)-sector eliminated tensors keep their old leading-order form. Any change in those sectors can only enter through higher-order feedback of the redesigned metric into the same lower-order coefficients, which is \(\mathcal O_{\ge 3}\) at the present expansion level.
\end{proof}

\begin{remark}
Proposition \ref{prop:blind} is the central structural point of the present note. A bridge term built only from \(\vartheta\) can cancel a convergence-generated longitudinal slot, but it cannot by itself eliminate an independent off-longitudinal eliminated tensor.
\end{remark}

\section{A Decisive Off-Longitudinal Surviving Remainder}

\begin{definition}[Rotationally nontrivial convergence-free redesigned regime]
\label{def:rot}
A reduced redesigned solution is said to lie in the \emph{rotationally nontrivial convergence-free regime} if
\begin{equation}
K\equiv 0,
\qquad
\Omega_{ij}^T\not\equiv 0,
\label{eq:rot}
\end{equation}
where
\[
\Omega_{ij}^T\equiv D_iW_j^T-D_jW_i^T
\]
is the transverse-vorticity two-form determined by the slice solve.
\end{definition}

\begin{remark}
Definition \ref{def:rot} is part of the physical scope already suggested by the mass-shaped gravity criterion: rotational dragging is one of the channels that the full bridge should eventually encode. The minimal convergence-response redesign did not remove or constrain that sector.
\end{remark}

\begin{proposition}[The redesigned bridge term vanishes on the convergence-free rotational regime]
\label{prop:convfree}
On the regime of Definition \ref{def:rot}, the bridge-redesign slot vanishes at quadratic order:
\begin{equation}
\Phi_\vartheta\equiv 0,
\qquad
\PiTheta_{AB}=0
\quad\text{to quadratic order}.
\label{eq:convfree}
\end{equation}
\end{proposition}

\begin{proof}
If \(K\equiv0\), then the leading convergence scalar \(\vartheta\) also vanishes at the same order used by the redesign note, so its elliptic potential \(\Phi_\vartheta\) vanishes there as well. Equation \eqref{eq:pitheta} then gives \eqref{eq:convfree}.
\end{proof}

\begin{proposition}[Leading vorticity effective action and tensor]
\label{prop:W}
On the regime of Definition \ref{def:rot}, the leading redesigned vorticity effective action is
\begin{equation}
\Gamma_{W,\mathrm{eff}}^{\mathrm{br},(2)}
=
-\frac{\kappa_\Omega}{4}
\int_{\mathbb R^{1+3}}\Omega_{ij}^T\Omega^{T,ij}\,dt\,d^3x,
\label{eq:GammaW}
\end{equation}
and its metric variation gives
\begin{equation}
T_{\mu\nu}^{(W,\mathrm{br,eff},2)}
=
\kappa_\Omega
\left(
\Omega_{\mu\alpha}^T\Omega_\nu^{T\,\alpha}
-\frac14 g_{\mu\nu}^{\mathrm{br}}\Omega_{\alpha\beta}^T\Omega^{T,\alpha\beta}
\right).
\label{eq:TW}
\end{equation}
Its Hamiltonian projection is
\begin{equation}
\rho_W^{\mathrm{br},(2)}
\equiv
n^\mu n^\nu T_{\mu\nu}^{(W,\mathrm{br,eff},2)}
=
\frac{\kappa_\Omega}{4}\abs{\Omega^T}^2.
\label{eq:rhoW}
\end{equation}
\end{proposition}

\begin{proof}
Once the longitudinal convergence source is set to zero at the relevant order, the leading off-longitudinal auxiliary functional comes from the unchanged transverse-vorticity block
\[
-\frac{\kappa_\Omega}{4}\Omega_{AB}\Omega^{AB}
\]
already present in the reduced quotient action. Restricting that block to the spatially transverse vorticity solve on the redesigned branch gives \eqref{eq:GammaW}. Let \(n^\mu\) be the future unit normal of the unit-lapse quotient slicing, and extend the spatial two-form \(\Omega_{ij}^T\) to a spacetime two-form \(\Omega_{\mu\nu}^T\) by requiring \(n^\mu\Omega_{\mu\nu}^T=0\). Varying the quadratic two-form energy with respect to the operational metric then produces the standard Maxwell-type stress tensor \eqref{eq:TW}. Contracting with the unit normal gives \eqref{eq:rhoW}.
\end{proof}

\begin{theorem}[Redesigned bridge tensor verdict]
\label{thm:verdict}
On the current redesigned bridge branch:
\begin{enumerate}
    \item the actual eliminated observer equation is
    \begin{equation}
    G_{\mu\nu}[g^{\mathrm{br}}]
    =
    8\pi G_N\,T_{\mu\nu}^{\mathrm{matter}}
    +T_{\mu\nu}^{(Q,\mathrm{br,eff})}
    +T_{\mu\nu}^{(W,\mathrm{br,eff})}
    +\ReLong;
    \label{eq:full}
    \end{equation}
    \item \(\ReLong\) is only a higher-order longitudinal residual, not a proof of all-order cancellation;
    \item the off-longitudinal remainder survives genuinely: in the rotationally nontrivial convergence-free regime of Definition \ref{def:rot},
    \begin{equation}
    \ReEff
    =
    T_{\mu\nu}^{(W,\mathrm{br,eff},2)}
    +\mathcal O_{\ge 3},
    \label{eq:survive}
    \end{equation}
    with nonzero Hamiltonian density \eqref{eq:rhoW};
    \item therefore the minimal convergence-response bridge does \emph{not} close the observer tensor. It removes the old same-scope longitudinal slot but only postpones the bridge problem.
\end{enumerate}
\end{theorem}

\begin{proof}
Item (1) is Proposition \ref{prop:decomp}. Item (2) is Proposition \ref{prop:longcancel}. For item (3), combine Proposition \ref{prop:convfree} with Proposition \ref{prop:blind} and Proposition \ref{prop:W}. On the regime \(K\equiv0\), the redesigned bridge slot itself vanishes at the same order, so the quadratic remainder comes from the transverse-vorticity effective tensor. Since \(\Omega_{ij}^T\not\equiv0\), equation \eqref{eq:rhoW} shows that this tensor has a nonzero observer-side Hamiltonian projection. Such a nonzero Hamiltonian density produced directly by metric variation of the effective action is not removed by the vanishing slice equations and is not merely a gauge artifact. Hence \(\ReEff\not\equiv0\) on that regime.

Item (4) follows immediately. Exact bridge closure would require the right-hand side of \eqref{eq:full} to vanish identically or to collapse to a pure gauge/constraint-exact expression on the full redesigned branch. The quadratic vorticity tensor \eqref{eq:TW} already rules that out.
\end{proof}

\begin{corollary}[Status of the higher-order longitudinal remainder]
\label{cor:high}
Even apart from the off-longitudinal obstruction of Theorem \ref{thm:verdict}, the redesigned bridge leaves an unresolved higher-order longitudinal remainder \(\ReLong=\mathcal O_{\ge 3}\). The same-scope redesign therefore does not provide an all-order longitudinal identity.
\end{corollary}

\begin{proof}
This is exactly Proposition \ref{prop:longcancel}.
\end{proof}

\begin{corollary}[Minimal bridge-first redesign is only a postponement]
\label{cor:postpone}
The minimal convergence-response redesign is a positive same-scope repair, not a full redesigned-bridge closure. It postpones the original longitudinal obstruction to higher order while leaving a genuine off-longitudinal observer remainder.
\end{corollary}

\begin{proof}
Combine Theorem \ref{thm:verdict} with Corollary \ref{cor:high}.
\end{proof}

\section{What the Verdict Changes}

\begin{proposition}[The next redesign may now move beyond bridge metric alone]
\label{prop:next}
Because the full redesigned-bridge closure test is now negative, it is finally honest to open the second redesign note that reconsiders the single-metric matter route.
\end{proposition}

\begin{proof}
The previous redesign note imposed the correct discipline: first test whether the smallest bridge-only change removes the observer remainder before reopening the matter route. Theorem \ref{thm:verdict} shows that this bridge-only test fails at full eliminated-tensor scope. Therefore the recorded conditional trigger for a second redesign note has now been reached.
\end{proof}

\begin{remark}
Proposition \ref{prop:next} does not claim that every future successful route must abandon the bridge-metric line completely. It records only the sequence verdict requested here: after the full redesigned-bridge test remains negative, a disciplined matter-route reconsideration is now honest.
\end{remark}

\section{Conclusion}

The present note performs the first full redesigned-bridge closure test that the previous same-scope redesign note deliberately left open.
\begin{enumerate}
    \item It keeps the positive quotient reduced-system/theorem/decay/asymptotic line closed and keeps the single-metric matter route fixed throughout the test.
    \item It derives the actual eliminated observer equation on the revised bridge metric.
    \item It shows that the old quadratic longitudinal quotient tensor is cancelled, but only at that same scope, leaving a higher-order longitudinal residual \(\ReLong=\mathcal O_{\ge 3}\).
    \item It shows that the redesigned bridge is blind at leading order to the off-longitudinal \(Q\)- and vorticity-sector effective tensors.
    \item It proves a decisive negative verdict on a convergence-free rotational regime: the transverse-vorticity effective tensor survives with positive observer-side Hamiltonian density, so the redesigned observer tensor is still not zero.
\end{enumerate}

The minimal convergence-response bridge therefore does not achieve full observer-tensor closure. It was the right bridge-first test, and it was worth doing, but its success was only same-scope and only longitudinal. The full redesigned-bridge verdict is that the bridge problem is postponed rather than solved. With that verdict now recorded, the next honest manuscript may finally reconsider the single-metric matter route itself.

\end{document}
